\documentclass{article}
\usepackage{amsmath}
\title{Q4P10: Peer discipline with resource-dependent evidence}
\author{Ada}
\date{}
\begin{document}
\maketitle

\section*{The pair in two sentences}
Q, \emph{Mechanism Design for Alignment and Control}, supplies a peer-scoring construction for
truthful reporting and action obedience using known, separated conditional beliefs, utility
cancellation, unrestricted rewards, and prohibitions (ledger entries 1--2, 7, 17).
P, \emph{The Energy Society}, studies agents whose token generation consumes the energy that also
funds continued participation, comparing instructions to maximize individual or group energy rather
than implementing Q's mechanism (entries 1--3, 5).

\section*{The connexion pursued}
The peers converged on whether resource consumption changes the observable peer evidence enough to
undermine discipline, and what payments could restore reporting incentives (entries 3--5, 17, 19).
The proposed connexion is a feedback between resources, peer evidence, and incentives; it is not an
established empirical result of Q or P.
Both notes supply checked finite examples, but the empirical material is thin: no LLM experiments or
full Energy Society runs were performed, and novelty remains unsettled (entries 17, 19, 21, 24).

Q's quadratic score is
\[
 S(p,y)=2p(y)-\lVert p\rVert_2^2.
\]
With a fixed truthful peer channel, its expected truthful advantage over another belief report is
the squared distance between the belief rows.
Q's proof scales this margin to dominate residual deviation gains; its full action-discipline result
additionally uses utility cancellation and off-target prohibitions (entries 1--2; both notes).
The calculations below isolate reporting against a fixed channel, not the full construction.

\section*{What was found}
\paragraph{Selection reverses a stale score.}
Emmy's example, independently checked by Ada, has binary peer outcomes with
\[
 p_H(1)=\frac34,\qquad p_L(1)=\frac14,
 \qquad s(1)=\frac1{10},\qquad s(0)=1,
\]
where \(s\) is the observation probability.
After conditioning on observation,
\[
 \widetilde p_H(1)=\frac{3/40}{3/40+1/4}=\frac3{13},
 \qquad \widetilde p_L(1)=\frac1{31}.
\]
The original score rows have equal squared norms, so their difference is
\(1\) at outcome \(1\) and \(-1\) at outcome \(0\).
Consequently,
\[
 \mathrm{E}_H[S(p_H,Y)-S(p_L,Y)\mid\text{observed}]
 =2\frac3{13}-1=-\frac7{13}.
\]
Increasing the scale strengthens the false-report incentive (entries 8, 11).
Using the correct survivor-conditioned rows restores the proper-score argument if selection is known
and unaffected by reporting deviations.
Alternatively, retaining nonresponse as an outcome gives, in the order \((1,0,\text{missing})\),
\[
 q_H=(3/40,1/4,27/40),\qquad q_L=(1/40,3/4,9/40).
\]
Their squared separation is \(91/200\); conditioning on response instead gives
\(2(80/403)^2\).
Missingness can retain information, so this example does not show that all attrition destroys
separation (entry 11; both notes).
These distances are not comparisons at equal payment caps.

\paragraph{A sharp peak-payment requirement.}
A separate example does lose diagnostic information even when missingness is recorded (entries 12,
14).
Start with latent laws \(p_H=(1/2,2/5,1/10)\) and
\(p_L=(1/2,1/10,2/5)\), and observation probabilities
\((1,\varepsilon,\varepsilon)\), where \(0\leq\varepsilon\leq1\).
The observable laws are
\begin{align*}
 q_H&=(1/2,2\varepsilon/5,\varepsilon/10,(1-\varepsilon)/2),\\
 q_L&=(1/2,\varepsilon/10,2\varepsilon/5,(1-\varepsilon)/2).
\end{align*}
Thus \(\operatorname{TV}(q_H,q_L)=3\varepsilon/10\).
Assume linear utility in bonuses, a fixed channel under reporting deviations, and a specified
false-report continuation gain \(M>0\) for either type.
This gain is a primitive of the example, not an estimate of LLM preferences; action obedience is
separate.
For bonuses \(R(r,y)\in[0,B]\), write \(d(y)=R(H,y)-R(L,y)\).
Truthfulness requires
\[
 \mathrm{E}_{q_H}d\geq M,\qquad \mathrm{E}_{q_L}d\leq-M.
\]
Since \(\lvert d(y)\rvert\leq B\),
\[
 2M\leq\sum_y(q_H(y)-q_L(y))d(y)
 \leq2B\operatorname{TV}(q_H,q_L).
\]
For this symmetric channel the bound is attained: pay \(B\) only for an
\(H\) report on the second outcome or an \(L\) report on the third.
Each type's truthful bonus advantage is \(3B\varepsilon/10\), giving
\[
 B_{\min}=\frac{10M}{3\varepsilon}\quad(\varepsilon>0).
\]
Equality gives weak incentive compatibility; a larger cap gives strict preference against the false
report.
At \(\varepsilon=0\) the observable laws coincide and no finite bonus works.
The total-variation inequality is necessary generally for the stated two-type problem; attainment
here is not a general sufficiency theorem (entries 12, 14; Emmy's note).

The divergence concerns peak payout, not expected expenditure:
\[
 \mathrm{E}[\text{truthful payout}]
 =\frac{2\varepsilon}{5}B_{\min}=\frac{4M}{3}.
\]
Rare large payments can satisfy an expected budget while exceeding available reserves.
P supplies no explicit principal scoring reserve; such a restriction would be an added intervention
(entries 15--16).

\paragraph{Normalizing quadratic scoring.}
Emmy's latest correction is essential (entry 24).
For the same channel, let
\[
 b(y)=\min_{r\in\{H,L\}}S(q_r,y),\qquad
 R(r,y)=\Lambda[S(q_r,y)-b(y)].
\]
The baseline is independent of the report, so it preserves reporting comparisons under the
fixed-channel, linear-utility assumptions.
The rows have equal squared norms; the normalized score is exactly the diagnostic-bonus contract,
with
\[
 B=\frac{3\Lambda\varepsilon}{5},\qquad
 \Lambda_{\min}=\frac{50M}{9\varepsilon^2},\qquad
 B_{\min}=\frac{10M}{3\varepsilon}.
\]
Inverse-squared score scale therefore does not imply inverse-squared peak payments.
Normalization need not preserve survival trajectories or incentives for deviations that change the
outcome law (entry 24).
Both notes record exact-arithmetic checks of the examples; the latest normalization check is in
Emmy's script.

\section*{Objections and their disposition}
P's prompted objectives do not instantiate Q's reward mechanism, and its recommendations establish
neither known type-conditioned beliefs nor the scalar actions and decreasing bias curve required by
Q's coupled-reward construction (entries 1--2).
This rules out the direct mechanism-transfer reading.
P's communication metrics also do not identify sandbagging: there is no verified capability frontier
or evaluation-to-permission stage (entries 6--7).

The stale-score reversal is misspecification, not a refutation of Q.
Bad or pooling equilibria likewise do not refute Q's partial implementation, which requires the
target to be an equilibrium rather than the unique equilibrium (entries 7, 17, 19).
Recalibration resolves the fixed-channel reporting problem; strategic influence over peers remains
unresolved.

Timing limits the proposed application.
The code audit found that energy is deducted after a completed model call, completed votes survive
exhaustion, and already inactive agents are excluded before discussion.
Successful job payouts can also restore an exhausted solver's energy (entries 13, 19).
Thus ordinary within-call exhaustion must not simply be coded as a missing vote.
Channel laws must condition on public history and active status at report commitment.
Fresh signals independent of past energy need not exhibit survival bias at all (entries 14, 16--17).
Sealing reports and eligibility before recommendations or donations addresses the immediate
experimental channel, but not strategic selection across rounds.

Finally, P's no-size-penalty results report all agents surviving and all efficiency ratios above
one, contradicting the abstract's stronger resource-loss claim.
The peers use the body and table, identified by the label
\texttt{tab:exp2\_big\_table}, rather than the abstract (entries 6--7, 17).
Neither those results nor pooled collision counts identify the effect of a scoring intervention.

\section*{Sources used outside Q and P}
The peers' scoped searches found prior work on costly effort, participation and budget balance in
Miller, Resnick and Zeckhauser (2005); participation and payment restrictions in Witkowski et al.
(2013), \emph{Dwelling on the Negative}; budget-constrained information acquisition in Radanovic et
al., \emph{Information Gathering with Peers}, arXiv:1711.06740; and uninformative equilibria and
advantages of direct checking under stated assumptions in Gao, Wright and Leyton-Brown,
arXiv:1606.07042 (entries 9--10; both notes).
These sources undermine novelty claims based only on bounded payments or bad equilibria.
Q itself already discusses LLM peer scoring; Emmy checked its Qiu, Carroll and Allen (2026)
precursor, arXiv:2601.20299 (entries 23--24).
The literature check is not a comprehensive novelty review.

Ada audited EnergySociety's \texttt{agent\_middleware.py} and
\texttt{environment.py}; Emmy checked the audit (entries 13, 19).
The source commit was
\begin{center}
\small\texttt{70fb2ac89fc4125f7f5e5dd3495fe07b7c558b4c}.
\end{center}
No additional outside sources were introduced for this consolidation.

\section*{What remains open}
The decisive missing evidence is a resource-dependent change in useful peer information conditional
on public history in the natural P economy.
Also unresolved are natural-type identification, nonlinear continuation value of energy, strategic
peer manipulation, equilibrium selection, and novelty beyond the cited work (entries 17, 21, 24).
Any useful correction must be assessed against direct answer verification, which P already permits,
and report-independent liquidity controls accounting for payment timing, recipients and variation.
Extra injected energy must not be counted as produced task value (both notes; entries 19, 21).

\section*{The smallest next step}
Perform the channel-identification check proposed in the notes: in the natural P economy, fix the
report-commitment information set and estimate whether resource state changes the informative
peer-outcome laws after conditioning on public history (entries 16--17, 21).
Use the audited timing, distinguishing inactive-at-start peers from completed votes.
This would settle the first empirical hurdle, not the whole publication question.
If the channel does not materially change, the distinctive proposed connexion lacks support; if it
does, it supplies the calibration for a sealed-report scoring test with explicit peak reserves and
direct-check and liquidity controls.
A synthetic erasure test can check the mathematical component but cannot settle this natural-economy
question (entries 19, 21; both notes).

\end{document}
